% ---
% Arquivo com a formatação do TCC.
% Este arquivo não deve ser modificado.
% ---
\usepackage[utf8]{inputenc}
\usepackage[T1]{fontenc}
\usepackage{amsmath}
\usepackage{amssymb,amsfonts,textcomp}
\usepackage{color}
\usepackage{array}
\usepackage{supertabular}
\usepackage{listings}         % Para as linguagens de programação
\usepackage{lastpage}		  % Usado pela Ficha catalográfica
\usepackage{indentfirst}	  % Indenta o primeiro parágrafo de cada seção.
\usepackage{hhline}
\usepackage{hyperref}
\usepackage[pdftex]{graphicx}
\graphicspath{ {./figuras/} }
\usepackage{python}
\usepackage{placeins}


% retira as mensagens de aviso do pacote Glossaries
\let\printglossary\relax
\let\theglossary\relax
\let\endtheglossary\relax
% coloca Seções em maiúsculas sem negrito no Sumário
\makeatletter
\let\oldcontentsline\contentsline
\def\contentsline#1#2{%
  \expandafter\ifx\csname l@#1\endcsname\l@section
    \expandafter\@firstoftwo
  \else
    \expandafter\@secondoftwo
  \fi
  {%
    \oldcontentsline{#1}{\normalfont\MakeTextUppercase{#2}}%
  }{%
    \oldcontentsline{#1}{#2}%
  }%
}
\makeatother
% ---
% Pacotes glossaries
% ---
\usepackage[subentrycounter,seeautonumberlist,nonumberlist=true]{glossaries}
% para usar o xindy ao invés do makeindex:
%\usepackage[xindy={language=portuguese},subentrycounter,seeautonumberlist,nonumberlist=true]{glossaries}
% ---
% Citações de referências no formato alfabético e negrito
\usepackage[alf, abnt-emphasize=bf]{abntex2cite} 
% Margens definidas em 25 mm para uso como documento em PDF
% Para imprimir use as seguintes margens:
% \usepackage[left=30mm, top=30mm, right=20 mm, bottom=20mm] {geometry}

\usepackage[margin=25 mm]{geometry}
%---
% O arquivo com o nome dos alunos e dos orientadores é lido aqui.
% Atualizar diretamente no arquivo
%---
%%%%%%%%%%%%%%%%%%%%%%%%%%%%%%%%%%%%%%%%%%%%%%%%%%%%%%%%%%%%%
% Definições de Macros utilizadas no TCC - ATUALIZE AQUI
%%%%%%%%%%%%%%%%%%%%%%%%%%%%%%%%%%%%%%%%%%%%%%%%%%%%%%%%%%%%%
\instituicao{UNIVERSIDADE FEDERAL DO RIO DE JANEIRO
\par
DEPARTAMENTO DE ENGENHARIA METALÚRGICA E DE MATERIAIS
\par
CURSO DE BACHARELADO EM ENGENHARIA METALÚRGICA}
% \title{MODELO DE MONOGRAFIA DE TRABALHO DE CONCLUSÃO DE CURSO\\ Modelo em Latex}
\title{CARACTERIZAÇÃO DA MICROESTRUTURA DE LIGAS DE ALTA ENTROPIA ATRAVÉS DE APRENDIZADO DE MÁQUINA}
\autor{Leonardo Pereira Rodrigues }
\orientador{Profa. Adriana da Cunha Rocha}
\coorientador{Profa. --}
\local{RIO DE JANEIRO, RJ - BRASIL}
\data{Junho de 2021}
\preambulo{Projeto de graduação submetido ao corpo docente do curso de engenharia metalúrgica da escola politécnica da Universidade Federal do Rio de Janeiro como parte dos requisitos necessários para a obtenção do grau de engenheiro metalúrgico.}




%%%%%%%%%%%%%%%%%%%%%%%%%%%%%%%%%%%%%%%%%%%%%%%%%%%%%%%%%%%
% Corrige a fonte dos capítulos, seções, resumos, etc.
%%%%%%%%%%%%%%%%%%%%%%%%%%%%%%%%%%%%%%%%%%%%%%%%%%%%%%%%%%%

\renewcommand{\ABNTEXchapterfont}{\bfseries \rmfamily}  % Capítulos em Bold e Maiúsculas
\renewcommand{\ABNTEXchapterfontsize}{\normalsize}
\renewcommand{\ABNTEXsectionfont}{\rmfamily}            %  Seções em  Maiúsculas apenas
\renewcommand{\ABNTEXsectionfontsize}{\normalsize}
\renewcommand{\ABNTEXsubsectionfont}{\bfseries}         % Subseções em Bold apenas
\renewcommand{\ABNTEXsubsectionfontsize}{\normalsize}
\renewcommand{\lstlistingname}{Código}                  % Nome para os códigos no texto
\renewcommand{\lstlistlistingname}{Lista de \lstlistingname s}
\makeatletter                                          % Configura a linha da lista de códigos
\renewcommand\l@lstlisting[2]{{\normalfont\@dottedtocline{1}{1.5em}{2em}{Código~#1}{#2}}}
\makeatother
\usepackage{url16023}  % para retirar < e > da URL nas referências.
%%%%%%%%%%%%%%%%%%%%%%%%%%%%%%%%%%%%%%%%%%%%%%%%%%%%%%%%%%%
% Criação de quadros com numeração
%%%%%%%%%%%%%%%%%%%%%%%%%%%%%%%%%%%%%%%%%%%%%%%%%%%%%%%%%%%
\newcommand{\quadroname}{Quadro}
\newcommand{\listofquadrosname}{Lista de quadros}

\newfloat[chapter]{quadro}{loq}{\quadroname}
\newlistof{listofquadros}{loq}{\listofquadrosname}
\newlistentry{quadro}{loq}{0}
%%%%%%%%%%%%%%%%%%%%%%%%%%%%%%%%%%%%%%%%%%%%%%%%%%%%%%%%%%%
% configurações para atender às regras da ABNT
%%%%%%%%%%%%%%%%%%%%%%%%%%%%%%%%%%%%%%%%%%%%%%%%%%%%%%%%%%%
\setfloatadjustment{quadro}{\centering}
\counterwithout{quadro}{chapter}
\renewcommand{\cftquadroname}{\quadroname\space} 
\renewcommand*{\cftquadroaftersnum}{\hfill--\hfill}
%---
% Configuração de posicionamento padrão
%---
\setfloatlocations{quadro}{hbtp}
%---
% Para ajudar nas tabelas e quadros
%---
\makeatletter
\newcommand\arraybslash{\let\\\@arraycr}
\makeatother
%%%%%%%%%%%%%%%%%%%%%%%%%%%%%%%%%%%%%%%%%%%%%%%%%%%%%%%
% Redefine a macro para imprimir a capa 
%%%%%%%%%%%%%%%%%%%%%%%%%%%%%%%%%%%%%%%%%%%%%%%%%%%%%%%
\renewcommand{\imprimircapa}{%
\begin{capa}%
\center
\imprimirinstituicao
\par
\vspace*{1cm}
\MakeUppercase{\imprimirautor}
\vfill
\begin{center}
\imprimirtitulo
\end{center}
\vfill
\imprimirlocal
\par
\imprimirdata
\vspace*{1cm}
\end{capa}
}
%%%%%%%%%%%%%%%%%%%%%%%%%%%%%%%%%%%%%%%%%%%%%%%%%%%%%%%
% Redefine a macro para imprimir a folho de rosto
%%%%%%%%%%%%%%%%%%%%%%%%%%%%%%%%%%%%%%%%%%%%%%%%%%%%%%%

\renewcommand{\imprimirfolhaderosto}{%
% \begin{figure}[b]
% \centering
\includegraphics{figuras/poli-ufrj-logo.png}
% \end{figure}
\begin{capa}%
\center
\par
\vspace*{1cm}
% \MakeUppercase{\imprimirautor}
\vfill
\begin{center}
\imprimirtitulo
\end{center}
\vfill
{\imprimirautor}
\vfill
\hspace*{\fill}\parbox[b]{.5\textwidth}{%
        \linespread{1.5}\selectfont
Projeto de graduação apresentado ao corpo docente do curso de engenharia metalúrgica da escola politécnica da Universidade Federal do Rio de Janeiro como parte dos requisitos necessários para a obtenção do grau de engenheiro metalúrgico.

}
\vfill
\flushright
Orientador: \imprimirorientador\\

Co-orientador: \imprimircoorientador\\
\vfill
\begin{center}
\large\imprimirlocal
\par
\large\imprimirdata
\vspace*{1cm}
\end{center}
\end{capa}
}
%%%%%%%%%%%%%%%%%%%%%%%%%%%%%%%%%%%%%%%%%%%%%%%%%%%%%%%
% Informações do PDF inseridas automaticamente
% Atualizar apenas as palavras-chave se necessário
% Não modificar as cores dos links e referências
%%%%%%%%%%%%%%%%%%%%%%%%%%%%%%%%%%%%%%%%%%%%%%%%%%%%%%%
\makeatletter
\hypersetup{
pdftitle={\@title},
pdfauthor={\@author},
pdfsubject={\imprimirpreambulo},
pdfkeywords={MONOGRAFIA}{CIÊNCIA DA COMPUTAÇÃO}{DCC-UFRJ},
pdfcreator={LaTeX with abnTeX2},
colorlinks=true,
linkcolor=black,
citecolor=black,
urlcolor=black
}
\makeatother
%%%%%%%%%%%%%%%%%%%%%%%%%%%%%%%%%%%%%%%%%%%%%%%%%%%%
% Definição das Linguagens de Programação
%%%%%%%%%%%%%%%%%%%%%%%%%%%%%%%%%%%%%%%%%%%%%%%%%%%%
\definecolor{dkgreen}{rgb}{0,0.6,0}
\definecolor{gray}{rgb}{0.5,0.5,0.5}
\definecolor{purple}{rgb}{0.8,0,0.3}
\definecolor{orange}{rgb}{1,0.4,0}
\definecolor{lightlightgray}{rgb}{.95,.95,.95}
\definecolor{lightgray}{rgb}{.9,.9,.9}
\definecolor{lightgray2}{rgb}{.85,.85,.85}
\definecolor{darkgray}{rgb}{.4,.4,.4}
%---
% Comando para inserir a listagem de código, tem 4 parâmetros
% Linguagem, Caption, Label, Nome do Arquivo
%---
\newcommand{\includecode}[4][C]{\mbox{\lstinputlisting[caption=#2, label=#3, escapechar=, style=custom#1]{#4}}}

%---
%Define características em comum e numeração sequencial 
%---
\lstset{numberbychapter=false,framexleftmargin=5mm,  frame=shadowbox, rulesepcolor=\color{gray}}
%---
% A leitura da formatação das linguagens é feita aqui. 
% Modifique diretamente no arquivo.
%
%%%%%%%%%%%%%%%%%%%%%%%%%%%%%%%%%%%%%%%%%%%%%%%%%%%%
% Aqui você pode personalizar a linguagem C  
%%%%%%%%%%%%%%%%%%%%%%%%%%%%%%%%%%%%%%%%%%%%%%%%%%%%
\lstdefinestyle{customC}{
      language = C,
      breaklines=true,
      basicstyle=\footnotesize\ttfamily,
      keywordstyle=\bfseries\color{blue},
      commentstyle=\itshape\color{purple},
      identifierstyle=\color{black},
      stringstyle=\color{orange},
      showstringspaces=false}

% \includecode[Java]{Exemplo em Linguagem Java} {alg:codigo2}{codigos/codigo.java}

%%%%%%%%%%%%%%%%%%%%%%%%%%%%%%%%%%%%%%%%%%%%%%%%%%%%
% Aqui você pode personalizar a linguagem Java 
%%%%%%%%%%%%%%%%%%%%%%%%%%%%%%%%%%%%%%%%%%%%%%%%%%%%       
\lstdefinestyle{customJava}{
      language = Java,
      breaklines=true,
      basicstyle=\footnotesize\ttfamily,
      keywordstyle=\bfseries\color{blue},
      commentstyle=\itshape\color{purple},
      identifierstyle=\color{black},
      stringstyle=\color{orange},
      showstringspaces=false}
%%%%%%%%%%%%%%%%%%%%%%%%%%%%%%%%%%%%%%%%%%%%%%%%%%%%
% Aqui você pode personalizar a linguagem C++
%%%%%%%%%%%%%%%%%%%%%%%%%%%%%%%%%%%%%%%%%%%%%%%%%%%%  
\lstdefinestyle{customC++}{
      language = C++,
      breaklines=true,
      basicstyle=\footnotesize\ttfamily,
      keywordstyle=\bfseries\color{blue},
      commentstyle=\itshape\color{purple},
      identifierstyle=\color{black},
      stringstyle=\color{orange},
      showstringspaces=false}
%%%%%%%%%%%%%%%%%%%%%%%%%%%%%%%%%%%%%%%%%%%%%%%%%%%%
% Aqui você pode personalizar a linguagem Javascript
%%%%%%%%%%%%%%%%%%%%%%%%%%%%%%%%%%%%%%%%%%%%%%%%%%%%        
\lstdefinestyle{customJavaScript}{
      language = JavaScript,
      breaklines=true,
      basicstyle=\footnotesize\ttfamily,
      keywordstyle=\bfseries\color{blue},
      commentstyle=\itshape\color{purple},
      identifierstyle=\color{black},
      stringstyle=\color{orange},
      showstringspaces=false}
      
% Javascript
\lstdefinelanguage{JavaScript}{
    aboveskip=15pt,
    keywords={typeof, new, true, false, catch, function, return, null, catch, switch, var, if, in, while, do, else, case, break},
    keywordstyle=\color{blue}\bfseries,
    ndkeywords={class, export, boolean, throw, implements, import, this},
    ndkeywordstyle=\color{darkgray}\bfseries,
    identifierstyle=\color{black},
    sensitive=false,
    comment=[l]{//},
    morecomment=[s]{/*}{*/},
    commentstyle=\color{purple}\ttfamily,
    stringstyle=\color{red}\ttfamily,
    morestring=[b]',
    morestring=[b]",
    rulecolor=\color{lightgray2},
    breaklines=true,
    basicstyle=\footnotesize\ttfamily,
    frame=single,
    backgroundcolor=\color{lightlightgray}
}
%%%%%%%%%%%%%%%%%%%%%%%%%%%%%%%%%%%%%%%%%%%%%%%%%%%%
% Aqui você pode personalizar o formato JSON
%%%%%%%%%%%%%%%%%%%%%%%%%%%%%%%%%%%%%%%%%%%%%%%%%%%%        
\lstdefinestyle{customJSON}{
      language = JSON,
      breaklines=true,
      basicstyle=\footnotesize\ttfamily,
      keywordstyle=\bfseries\color{blue},
      commentstyle=\itshape\color{purple},
      identifierstyle=\color{black},
      stringstyle=\color{orange},
      showstringspaces=false}
      
% JSON
\lstdefinelanguage{JSON}{
    aboveskip=15pt,
    string=[s]{"}{"},
    stringstyle=\color{blue},
    comment=[l]{:},
    commentstyle=\color{black},
    rulecolor=\color{lightgray2},
    breaklines=true,
    basicstyle=\footnotesize\ttfamily,
    frame=single,
    backgroundcolor=\color{lightlightgray}
}
%%%%%%%%%%%%%%%%%%%%%%%%%%%%%%%%%%%%%%%%%%%%%%%%%%%%
% Aqui você pode personalizar a linguagem de montagem do SAPIENS
%%%%%%%%%%%%%%%%%%%%%%%%%%%%%%%%%%%%%%%%%%%%%%%%%%%%       
\lstdefinestyle{customSapiens}{
      language = Sapiens,
      breaklines=true,
      basicstyle=\footnotesize\ttfamily,
      keywordstyle=\bfseries\color{blue},
      commentstyle=\itshape\color{purple},
      identifierstyle=\color{black},
      stringstyle=\color{orange},
      showstringspaces=false}

% Linguagem de Montagem    
\lstdefinelanguage{Sapiens}{
    aboveskip=15pt,
    keywordstyle=\color{blue}\bfseries,
    comment=[l]{;},
    commentstyle=\color{purple},
    rulecolor=\color{lightgray2},
    breaklines=true,
    basicstyle=\footnotesize\ttfamily,
    frame=single,
    backgroundcolor=\color{lightlightgray},
    language=[x86masm]Assembler,
    morekeywords={DB, DS, DW, JN, JSR, LDA, ORG, STA, STS, TRAP}
}

%%%%%%%%%%%%%%%%%%%%%%%%%%%%%%%%%%%%%%%%%%%%%%%%%%%%
% Aqui você pode personalizar a linguagem Python 
%%%%%%%%%%%%%%%%%%%%%%%%%%%%%%%%%%%%%%%%%%%%%%%%%%%%       
\lstdefinestyle{customPython}{
      language = Python,
      breaklines=true,
      basicstyle=\footnotesize\ttfamily,
      keywordstyle=\bfseries\color{blue},
      commentstyle=\itshape\color{darkgray},
      identifierstyle=\color{black},
      stringstyle=\color{orange},
      showstringspaces=true}

      
      
% \lstset{language=Python, 
%         basicstyle=\ttfamily\small, 
%         keywordstyle=\color{keywords},
%         commentstyle=\color{comments},
%         stringstyle=\color{red},
%         showstringspaces=false,
%         identifierstyle=\color{green},
%         procnamekeys={def,class}}


%---
% FIM
%---

%=====================================================================
% Definição do conteúdo de "dedicatória"
%=====================================================================
\newcommand{\makededicationpage}{

    \vspace*{\fill}

  \hspace*{\fill}\parbox[b]{.5\textwidth}{%
  \begin{flushright}
\begin{dedicatoria}
\vspace*{\fill}
Dedico aos meus pais Eduardo e Sueli, e a todos que me deram força esperança de continuar até o fim.
\vspace*{\fill}
\end{dedicatoria}
\end{flushright}
  }
  
  
  \pagebreak
}{}

%=====================================================================